% ============================================================
% Resumen (abstract)
% ============================================================

\begin{abstract}
    \noindent En esta lista de problemas se estudian los fundamentos de las cadenas y los
    lenguajes formales sobre un alfabeto $\Sigma$. Se trabajó con los conjuntos
    $\Sigma^n$, $\Sigma^*$ y $\Sigma^+$, la longitud y la concatenación de cadenas, y
    operaciones entre lenguajes como la unión, la intersección, las potencias, la
    reflexión y las clausuras de Kleene y positiva. Mediante el Principio de
    Inducción Matemática se demostraron resultados como $|\Sigma^n| = |\Sigma|^n$,
    $|u^n| = n|u|$ y $(uv)^R = v^R u^R$, además de las propiedades de $A^*$ y $A^+$.
    Entre las observaciones principales destacan que la concatenación es asociativa
    pero no conmutativa, aunque $|uv| = |vu|$; que en general $L^* \neq L$; y que el
    lenguaje vacío $\varnothing$ debe distinguirse de $\{\lambda\}$. Estos resultados
    constituyen la base para el estudio posterior de los autómatas finitos, las
    expresiones regulares y las gramáticas formales.
\end{abstract}